% Placeholder tikz overview figure. Rendered inline by main.tex.
\begin{figure}[t]
  \centering
  \begin{tikzpicture}[
    node distance=1.1cm,
    box/.style={draw,rounded corners,align=center,minimum width=2.6cm,minimum height=0.9cm,font=\small},
    memory/.style={draw,rounded corners,align=center,minimum width=3.4cm,minimum height=1.2cm,font=\small,fill=blue!5},
    link/.style={->,>=latex,thick},
    softlink/.style={->,>=latex,dashed},
  ]
    \node[box] (turn) {\textbf{Turn}\\(\emph{session}, \emph{speaker}, \emph{text})};
    \node[box,right=1.1cm of turn] (writer) {\textbf{Writer}\\claim extraction};
    \node[memory,right=1.1cm of writer] (claim) {\textbf{Claim}\\$(s,p,o,I,\pi,v,\rho,\phi)$};
    \node[memory,below=0.9cm of claim,fill=green!5] (graph) {\textbf{Memory graph}\\support / contradict / supersede};
    \node[box,left=1.1cm of graph] (linker) {\textbf{Linker}\\joint LLM labeller};

    \draw[link] (turn) -- (writer);
    \draw[link] (writer) -- (claim);
    \draw[link] (claim) -- (graph);
    \draw[link] (linker) -- (graph);
    \draw[softlink] (claim) -- (linker);

    \node[box,below=0.9cm of graph,fill=orange!5] (retrieve) {\textbf{Truth-consistent retrieval}\\query parse $\to$ $k$-NN $\to$ max-consistent subgraph};
    \node[box,below=0.5cm of retrieve] (reader) {\textbf{Reader}\\(answer or abstain)};
    \draw[link] (graph) -- (retrieve);
    \draw[link] (retrieve) -- (reader);
  \end{tikzpicture}
  \caption{TTMG at a glance. At write time, the \emph{writer} turns each
    session into \emph{claim} nodes with validity intervals and provenance;
    the \emph{linker} labels candidate neighbours with one of
    \texttt{support}, \texttt{contradict}, \texttt{supersede}, or
    \texttt{unrelated}. At read time, a query parser extracts the temporal
    anchor, $k$-NN retrieves candidates from the \emph{active} subset, and
    the retriever returns the maximum consistent subgraph. The reader
    receives claims with their validity intervals and can be forced to
    abstain when two active claims disagree.}
  \label{fig:overview}
\end{figure}
